        	 	
        	 	\begin{rmk}\label{rmk:coord}
        	 		Suppose that we are in the setting of Case II, and consider $B_{\varrho_k}(p_k)\cap C_{\varrho_k}(q_k)$. Letting $\Phi_k : B_{\varrho_k}(p_k) \to B^{\R^{n+1}}_{\varrho_k/r_k}(0)$ and $\Psi_k : C_{\varrho_k}(q_k)\to C^{\R^{n+1}}_{\varrho_k/r_k}(-c)$ denote the blown-up normal and Fermi-coordinates respectively, one can easily check that 
        	 		$$\Phi_k\circ \Psi_k^{-1}:  \Psi_k(B_{\varrho_k}(p_k)\cap C_{\varrho_k}(q_k)) \longrightarrow \Phi_k(B_{\varrho_k}(p_k)\cap C_{\varrho_k}(q_k))$$
        	 		converges locally smoothly on $\Pi_1$ to the identity as $k\to \infty$. Therefore the resulting blow-up and conclusions in Case II can actually be taken exactly as in Case I with respect to normal coordinates centered at $p_k$, or indeed as stated in the Theorem. 
        	 	\end{rmk}
        	 	
        	 	\begin{rmk}\label{rmk:equivballs}
        	 			Referring to the statement above, we further note that in these coordinates geodesic balls in $N$, $B_r(p)$, are directly comparable to Euclidean balls: for any $T>0$ there exists a sequence $\left\{\beta_k(T)\right\}$ with $\beta_k(T) \searrow 1$ so that for all $p\in B_{Tr_k}(s_k)$, if we denote $\widetilde{p}\in C^{\R^{n+1}}_{T}(-c)$ the point corresponding to $p$ in these rescaled coordinates then, by abusing notation,  
        	 			$$B^{\R^{n+1}}_{r\beta_k^{-1}}(\widetilde{p}) \subset B_r(p)\subset B^{\R^{n+1}}_{r \beta_k}(\widetilde{p}).$$
        	 		 This equivalence will be exploited along the course of the following proof when transferring certain variational properties (typically: eigenvalue bounds) back and forth between geodesic balls and coordinate balls. As it is well-known, the same equivalence holds true in geodesic normal coordinates as well.
        	 	\end{rmk}	
        	 	
        
        	 	
        	 	
        	 	
        	 	
        	 	
        	 	
        	 	
        	 	
